% !PRINTED 357
\weberpage{357}
Diese Gleichung zeigt aber, dass die Curve $14^{\text{ten}}$ Grades:

\weq{19}{\chi(x) = 0,}

die von den $a$ gänzlich unabhängig ist, durch die Berührungspunkte der Doppeltangenten, aber durch keinen anderen Punkt der Curve $f(x)$ hindurchgeht.

Die Function $\chi(x)$ ist rational von den Coëfficienten der Gleichung $f(x)$ abhängig. Es giebt unendlich viele verschiedene solcher Curven, unter denen sich auch Covarianten von $f(x)$ finden. Man kann sie, wie \gesperrt{Hesse} nachgewiesen hat, in einfacher Weise durch Determinanten ausdrücken.

Hier ziehen wir daraus den Schluss:

2.\ \gesperrt{Eine Curve vierter Ordnung ohne singulären Punkt hat 28 Doppeltangenten.}

Einem Bedenken gegen diesen Schluss ist aber noch zu begegnen. Es wäre denkbar, dass die Curve $\chi$ die Curve $f$ berührt, oder dass die Curve $f$ durch singuläre Punkte der Curve $\chi$ hindurchgeht. Dann würde sich die Anzahl der Schnittpunkte und möglicherweise auch die Anzahl der Doppeltangenten vermindern, so dass unsere Schlussweise eigentlich nur lehrt, dass eine Curve vierter Ordnung nicht mehr als 28 Doppeltangenten haben kann. Dies Bedenken wird sich aber durch die folgenden Betrachtungen von selbst dadurch erledigen, dass wir Formen der Curvengleichung kennen lernen werden, bei denen die Existenz von 28 verschiedenen Doppeltangenten ersichtlich ist.

\begin{center}
\S.\,96.

\gesperrt{Die Steiner'schen Complexe.}
\end{center}

Wenn $x_1 = 0$ die Gleichung irgend einer Doppeltangente der Curve vierter Ordnung $f = 0$ ist, die wir kurz die Doppeltangente $x_1$ nennen, so muss, wenn man $x_1 = 0$ setzt, $f$ in ein Quadrat übergehen. Daraus ergiebt sich, dass $f$ von der Form sein muss:

\weq{1}{f = x_1 V - u^2,}

worin $V$ eine cubische, $u$ eine quadratische Form ist. Der Function $f$ kann aber auf dreifach unendlich viele Arten die Gestalt (1) gegeben werden. Denn bedeutet $p$ eine beliebige
%CONT>>
